problem:  
Let \( (M^n,g) \) be a complete, noncompact Riemannian manifold such that  

1. The sectional curvature satisfies \( \mathrm{sec}_M \ge 0 \) outside a compact subset \( K \subset M \);  
2. There exists a **unique tangent cone at infinity** \( C_\infty(M,p) = C(\Sigma) \), independent of the basepoint \(p\) (i.e., all pointed Gromov–Hausdorff rescalings \( (M, r_i^{-2} g, p) \to C(\Sigma) \) for \( r_i \to \infty \) converge to the same metric cone);  
3. The cone link \( \Sigma \) is a smooth Riemannian manifold of dimension \( n-1 \).  

**Question:**  
Does there exist a compact, connected, smooth, totally convex submanifold \( S \subset M \) (a “generalized soul”) such that  

- \( M \) is diffeomorphic to the total space of the normal bundle \( \nu(S) \), and  
- the homeomorphism type of \( S \) is uniquely determined by the homeomorphism type of the cross section \( \Sigma \) of the tangent cone at infinity?  

More strongly, does the uniqueness of the tangent cone at infinity imply a corresponding **topological rigidity** of the soul—namely, that any two manifolds satisfying the same asymptotic conditions and sharing the same cone link \( \Sigma \) are homeomorphic (or even diffeomorphic) outside compact subsets?

Why is it a "good" problem:  
This problem probes the uncharted boundary between **global topology** and **asymptotic geometry**. It proposes a conjectural extension of the Cheeger–Gromoll Soul Theorem from global nonnegative curvature to the case where curvature is nonnegative only outside a compact set—but augmented by strong asymptotic regularity in the form of a unique tangent cone at infinity.  

The question asks whether this asymptotic cone, a purely **metric** object describing the “shape at infinity,” controls the **smooth differential-topological core** of the manifold. If true, it would establish a deep rigidity principle: that large-scale geometric behavior determines compact topological structure. Addressing it would require synthesizing techniques from **Cheeger–Colding theory**, **Alexandrov geometry**, and **Riemannian topology**, uniting synthetic metric geometry with differential-geometric analysis.  

The statement is elementary to express yet potentially foundational—its resolution could inspire a new “asymptotic soul theory” illuminating how the geometry near infinity dictates the global topology of noncompact spaces. This combination of simplicity, depth, and cross-field significance makes it an excellent research-level problem.